\scanpage{143}
% archive foliation: 142
$\cdot$ Une courbe elliptique $E$ sur $\mathbb C$ est déterminée par
\begin{equation}
\left\{
\begin{array}{l}
\text{a) ``réseau'' }H,\ \deleted{\unsure{\phi}}\ \mathbb Z\text{-module libre de rang }2
\qquad\deleted{\unsure{(H=1}}\\
(H=\pi_1(E,0_E)\simeq H_1(E,\mathbb Z)\simeq H^1(E,\mathbb Z)^\vee)\\
\text{b) Une structure complexe sur }V=H\otimes_{\mathbb Z}\mathbb R\ \deleted{\unsure{\simeq V}}\
(\simeq t_E,\\
\qquad\text{espace tangent à l'origine de }E).
\end{array}
\right.
\tag{1}
\end{equation}

La donnée a) permet de reconstituer $E$ comme groupe de Lie réel par
\begin{equation}
E\simeq H\otimes_{\mathbb Z}\mathbb R/H
\tag{2}
\end{equation}
et la donnée b) redonne la structure complexe sur $E$, par passage au quotient.

Quand on est sur une base $S$, espace $\mathbb C$-analytique, la donnée a) correspond à un système local $\mathcal H$ sur $S$ --- faisceau en $\mathbb Z$-modules loc libres de rang $2$ sur $S$, d'où un fibré vectoriel $\mathbb R$-analytique $\mathcal H\otimes_{\mathbb Z}\mathbb R$ sur $S$. La structure b) correspond à une structure complexe \deleted{sur} de ce dernier, en tant que fibré vectoriel \inserted{complexe} de rang $1$, de façon que les sections locales de $\mathcal H$ deviennent des sections analytiques complexes. Si \deleted{$\unsure{e}$} $z_1,z_2$ sont deux sections de $\mathcal H$ formant
